\section*{Chapitre \ref{ch:g3:continents-and-oceans} --- Le globe~: continents et océans}

\begin{solution}{exo:g3:continents-and-oceans:1}
L'Asie, l'Afrique, l'Amérique du Nord, l'Amérique du Sud, l'Antarctique,
l'Europe et l'Australie. Le plus grand est l'Asie~; le plus petit est
l'Australie.
\end{solution}

\begin{solution}{exo:g3:continents-and-oceans:2}
Les \omterm{def:g2:water-on-earth:ocean}{océans} Pacifique, Atlantique, Indien, Austral et Arctique. Le plus
grand est l'\omterm{def:g2:water-on-earth:ocean}{océan} Pacifique~; le plus petit est l'\omterm{def:g2:water-on-earth:ocean}{océan} Arctique.
\end{solution}

\begin{solution}{exo:g3:continents-and-oceans:3}
Le \omterm{def:g3:continents-and-oceans:poles}{pôle Nord} est en haut d'un \omterm{def:g3:continents-and-oceans:globe}{globe}~; le \omterm{def:g3:continents-and-oceans:poles}{pôle Sud} est en bas~; l'\omterm{def:g3:continents-and-oceans:equator}{équateur}
fait tout le tour du milieu.
\end{solution}

\begin{solution}{exo:g3:continents-and-oceans:4}
L'Europe~: \omterm{def:g3:continents-and-oceans:equator}{hémisphère nord}. L'Australie~: \omterm{def:g3:continents-and-oceans:equator}{hémisphère sud}. L'Antarctique~:
\omterm{def:g3:continents-and-oceans:equator}{hémisphère sud}. L'Amérique du Nord~: \omterm{def:g3:continents-and-oceans:equator}{hémisphère nord}.
\end{solution}

\begin{solution}{exo:g3:continents-and-oceans:5}
L'\omterm{def:g2:water-on-earth:ocean}{océan} Arctique entoure le \omterm{def:g3:continents-and-oceans:poles}{pôle Nord}. L'\omterm{def:g2:water-on-earth:ocean}{océan} Austral entoure
l'Antarctique.
\end{solution}

\begin{solution}{exo:g3:continents-and-oceans:6}
De l'Afrique à l'Amérique du Sud~: l'\omterm{def:g2:water-on-earth:ocean}{océan} Atlantique. De l'Australie à
l'Amérique du Sud, vers l'est~: l'\omterm{def:g2:water-on-earth:ocean}{océan} Pacifique. De l'Afrique à
l'Australie~: l'\omterm{def:g2:water-on-earth:ocean}{océan} Indien.
\end{solution}

\begin{solution}{exo:g3:continents-and-oceans:7}
Sam vit près de l'\omterm{def:g3:continents-and-oceans:equator}{équateur}, où il fait chaud toute l'année. Bo vit près
d'un \omterm{def:g3:continents-and-oceans:poles}{pôle}, où il fait froid toute l'année.
\end{solution}

\begin{solution}{exo:g3:continents-and-oceans:8}
$10 - 7 = 3$~: son doigt tombe sur la terre 3 fois. C'est bien ce que l'on
peut attendre~: environ 71 carrés sur 100 de la Terre sont de l'eau, donc
environ 7 fois sur 10 le doigt s'arrête sur la \omterm{def:g2:water-on-earth:ocean}{mer}. (Avec seulement
10 lancers, elle pourrait aussi obtenir 6 ou 8~: le compte change d'une
partie à l'autre.)
\end{solution}

\begin{solution}{exo:g3:continents-and-oceans:9}
Pour dessiner le \omterm{def:g3:continents-and-oceans:world-map}{planisphère}, on a coupé le \omterm{def:g3:continents-and-oceans:globe}{globe} le long d'une ligne
passant par le milieu de l'\omterm{def:g2:water-on-earth:ocean}{océan} Pacifique, puis on l'a mis à plat. Une
moitié du Pacifique se retrouve au bord gauche de la \omterm{def:g2:drawing-to-map:map}{carte}, l'autre moitié
au bord droit. Il n'y a qu'un seul \omterm{def:g2:water-on-earth:ocean}{océan} Pacifique sur la vraie Terre~: les
deux bords de la \omterm{def:g2:drawing-to-map:map}{carte} sont la même ligne sur le \omterm{def:g3:continents-and-oceans:globe}{globe}.
\end{solution}

\begin{solution}{exo:g3:continents-and-oceans:10}
Tom a tort. Du plus grand au plus petit, les \omterm{def:g3:continents-and-oceans:continent}{continents} sont l'Asie,
l'Afrique, l'Amérique du Nord, l'Amérique du Sud, l'Antarctique, l'Europe
et l'Australie~: l'Antarctique n'est que le cinquième, bien plus petit que
l'Asie. Le \omterm{def:g3:continents-and-oceans:world-map}{planisphère} le trompe parce qu'une \omterm{def:g2:drawing-to-map:map}{carte} plate étire beaucoup
les terres près des \omterm{def:g3:continents-and-oceans:poles}{pôles}. L'Antarctique, qui entoure le \omterm{def:g3:continents-and-oceans:poles}{pôle Sud}, est
étiré en une longue bande tout le long du bas de la \omterm{def:g2:drawing-to-map:map}{carte}. Sur le \omterm{def:g3:continents-and-oceans:globe}{globe},
c'est une masse de terre ronde, bien plus petite que l'Asie.
\end{solution}
